% draft_data.tex
% Style: Academic Paper
% Target: 2--3 pages (~1,200--1,600 words of body text)

\section{Data}
\label{sec:data}

\subsection{Source and Sample Construction}
\label{sec:data_source}

All data are sourced from the Bloomberg Terminal via the Bloomberg Desktop API (\texttt{blpapi}).
The sample universe consists of firms in the S\&P~500 index as of the initial pull date, yielding
503 unique equity securities identified by their Bloomberg ticker (e.g., \texttt{AAPL US Equity}).
I construct two panel datasets:

\begin{enumerate}
    \item \textbf{Primary panel.} Daily observations from January~1, 2025 through March~25, 2026,
    comprising 160,487 firm-day observations. This panel contains the full set of treatment,
    outcome, and control variables used in the FE-OLS and DoubleML specifications.

    \item \textbf{Extended panel.} Daily observations from January~1, 2016 through February~27,
    2026, comprising approximately 1.19~million firm-day observations. This panel contains a
    reduced variable set (prices, volume, market capitalization, sentiment scores) and is used
    exclusively for the temporal analysis that replicates \citet{gu2020} across rolling windows.
\end{enumerate}

Non-trading days (weekends, holidays) are excluded by dropping observations where both opening
and closing prices are missing. Bloomberg-specific missing value codes (\texttt{\#N/A N/A},
\texttt{\#N/A Field Not Applicable}) are recoded as \texttt{NA} prior to any numerical
transformation.

\subsection{Variable Definitions}
\label{sec:data_variables}

Table~\ref{tab:variables} defines all variables used in the analysis. Variables fall into four
categories: outcome, treatment, controls, and derived.

\begin{table}[htbp]
\centering
\caption{Variable Definitions}
\label{tab:variables}
\small
\begin{tabular}{lll}
\hline\hline
Variable & Bloomberg Field & Definition \\
\hline
\multicolumn{3}{l}{\textit{Outcome}} \\
\texttt{return} & --- & $\text{px\_close}_{it} - \text{px\_open}_{it}$ (open-to-close price change, \$) \\
\addlinespace[0.5ex]
\multicolumn{3}{l}{\textit{Treatment}} \\
\texttt{twitter\_sent}       & \texttt{TWITTER\_SENTIMENT\_DAILY\_AVG}  & Firm-day Twitter sentiment score ($[-1, 1]$) \\
\texttt{twitter\_neg\_count} & \texttt{TWITTER\_NEG\_SENTIMENT\_COUNT}  & Count of negative-polarity tweets \\
\texttt{news\_sent}          & \texttt{NEWS\_SENTIMENT\_DAILY\_AVG}     & Firm-day news sentiment score ($[-1, 1]$) \\
\addlinespace[0.5ex]
\multicolumn{3}{l}{\textit{Controls}} \\
\texttt{px\_high}        & \texttt{PX\_HIGH}              & Intraday high price (\$) \\
\texttt{px\_low}         & \texttt{PX\_LOW}               & Intraday low price (\$) \\
\texttt{mkt\_cap}        & \texttt{CUR\_MKT\_CAP}         & Market capitalization (millions, \$) \\
\texttt{total\_equity}   & \texttt{TOTAL\_EQUITY}         & Book equity (millions, \$) \\
\texttt{debt\_to\_equity} & \texttt{TOT\_DEBT\_TO\_TOT\_EQY} & Total debt / total equity (\%) \\
\texttt{volume}          & \texttt{PX\_VOLUME}            & Daily share volume (shares) \\
\texttt{rsi\_30}         & \texttt{RSI\_30D}              & 30-day Relative Strength Index \\
\texttt{ma\_50}          & \texttt{MOV\_AVG\_50D}         & 50-day moving average of price (\$) \\
\addlinespace[0.5ex]
\multicolumn{3}{l}{\textit{Derived}} \\
\texttt{lag$k$}          & ---                            & $\text{return}_{i,t-k}$ for $k \in \{1,2,3,5,7\}$ \\
\texttt{bid\_ask\_spread} & \texttt{AVERAGE\_BID\_ASK\_SPREAD} & Daily average bid-ask spread (\$) \\
\hline\hline
\end{tabular}
\end{table}

The dependent variable, \texttt{return}, is the raw dollar price change from market open to
market close on day $t$ for firm $i$. This differs from the percentage return used in some prior
studies \citep{gu2020}; the dollar denomination preserves interpretability of treatment
coefficients as dollar-per-unit sentiment effects, and the firm fixed effects absorb
cross-sectional differences in share price levels.

The primary treatment variable, \texttt{twitter\_sent}, is Bloomberg's proprietary daily average
Twitter sentiment score. Bloomberg applies natural language processing to the universe of public
tweets mentioning a given firm's name, ticker, or cashtag, assigning each tweet a score on $[-1,
+1]$ and averaging across all tweets within the calendar day. The variable
\texttt{twitter\_neg\_count} counts only those tweets classified with negative polarity, following
the construction of \citet{teti2019}. The news sentiment variable, \texttt{news\_sent}, is
constructed analogously from Bloomberg-sourced news articles.

\subsection{Measurement Limitations}
\label{sec:data_limitations}

Three features of the Bloomberg sentiment data warrant discussion.

\textbf{Proprietary classification.} Bloomberg does not disclose the dictionary, model
architecture, or training data underlying its sentiment classification. This precludes replication
with alternative NLP tools and prevents diagnostic checks on whether specific tweet categories
(e.g., bot-generated, reply threads, retweets) are included or filtered. The
\citet{loughran2011} critique of general-purpose sentiment dictionaries applies with unknown
force: it is possible that Bloomberg's classifier misattributes domain-specific financial language.

\textbf{Temporal aggregation.} The sentiment scores are daily averages. Bloomberg aggregates all
tweets within a 24-hour window without distinguishing pre-market, intraday, and after-hours
activity. Because the outcome variable is measured from market open to market close, a tweet
posted at 9:00~PM after market close is counted in the same day's sentiment as one posted at
10:00~AM. This aggregation mismatch biases the design against detecting effects from tweets that
post-date the close, and may introduce reverse causality if after-hours tweets respond to the
day's price action that has already been recorded in \texttt{return}.

\textbf{Selection in sentiment coverage.} Not all firm-days have Twitter sentiment scores.
Smaller firms, firms with less public visibility, or firms experiencing quiet periods may have
fewer or no tweets classified. Missing sentiment days are not randomly distributed across firms or
time: large-cap, high-volume firms in the technology and consumer discretionary sectors receive
disproportionate coverage. This selection pattern means that the estimated treatment effect is
identified primarily from firms with active social media presence and may not generalize to the
full S\&P~500 universe. Summary statistics in Table~\ref{tab:sumstats} document the extent of
missingness across key variables.

\subsection{Sample Properties}
\label{sec:data_sample}

The primary panel spans 503 firms over approximately 320 trading days, yielding a balanced panel
with minor gaps due to delistings, IPOs within the sample window, and missing Bloomberg fields.
The effective sample size for the full FE-OLS specification (Model~2) drops to 147,128
observations after listwise deletion of rows with missing controls, primarily
\texttt{total\_equity} and \texttt{debt\_to\_equity} for newly listed firms. The baseline
specification (Model~1, sentiment only) retains the full 160,103 observations.

For the DoubleML specifications, XGBoost handles missing covariates natively---observations are
not dropped for missing controls. However, observations with missing treatment
(\texttt{twitter\_sent} or \texttt{news\_sent}) are necessarily excluded. The effective sample
size for DoubleML is reported in each results table.

The extended panel for temporal analysis covers 2,633 unique trading days across the same 503
firms. Coverage of Bloomberg Twitter sentiment begins in earnest around 2017 for most firms,
with sporadic coverage in 2016. The Fama-MacBeth regressions in the temporal analysis use only
dates with at least 50 firms reporting non-missing sentiment, following the minimum
cross-section requirement of \citet{gu2020}.

% Pointer to summary statistics table
% \input{output/sumstats_table.tex}


% =========================================================
% MISSING / EXPAND NOTES
% =========================================================
% [MISSING: Summary statistics table (Table~\ref{tab:sumstats}) showing mean,
%   SD, min, p25, median, p75, max, N for all key variables. Generate from
%   panel_long.csv and add as output/sumstats_table.tex]
% [MISSING: Correlation matrix of treatment and control variables, or at
%   minimum the correlation between twitter_sent and news_sent]
% [EXPAND: Sector composition of the 503 firms — GICS sector breakdown to
%   characterize which industries drive the panel]
% [EXPAND: Time-series plot of daily cross-sectional mean twitter_sent to
%   visualize sentiment dynamics over the sample period]
